\section{Repaso del paradigma ReAct}

\subsection{De CoT a ReAct}

ReAct es uno de los trabajos más emblemáticos en el campo de los agentes (equipo de Yao Shunyu (姚顺雨), publicado a finales de 2022 / principios de 2023 en NeurIPS, con más de 4,000 citas). Su línea de evolución es la siguiente:

\begin{enumerate}
    \item \textbf{CoT (Chain of Thought)}: orientado a responder preguntas: ``Let's think step by step'', que inspiró el Prompt Engineering posterior.
    \item \textbf{Act Only}: el modelo de lenguaje puede interactuar con el entorno, produciendo una acción (Action) y recibiendo del entorno una observación (Observation), pero sin un proceso de pensamiento explícito.
    \item \textbf{ReAct = Reasoning + Acting}: antes de producir la acción se realiza un razonamiento (Thought), lo que reduce la alucinación del modelo y logra un mejor grounding hacia acciones razonables.
\end{enumerate}

\begin{importantbox}{El ciclo central de ReAct}
Thought $\rightarrow$ Action $\rightarrow$ Observation $\rightarrow$ Thought $\rightarrow$ ... $\rightarrow$ Final Answer

ReAct inspiró el diseño y desarrollo de los agentes posteriores, y posiblemente también inspiró la aparición de la característica de Function Calling en la API de GPT (junio de 2023).
\end{importantbox}

\begin{knowledgebox}{Sobre los Reasoning Model}
Los Reasoning Model modernos ya han aprendido a pensar antes de responder una pregunta, sin necesidad de indicarles explícitamente mediante el prompt que «primero deben pensar». Más adelante, al presentar los Reasoning Model, se profundizará en sus normas y estándares de uso.
\end{knowledgebox}

\subsection{Resumen de la sección}
ReAct combina el razonamiento (Reasoning) con la acción (Acting), y constituye la base del diseño de agentes. Entender su tríada Thought-Action-Observation es el requisito previo para entender todos los marcos de agentes posteriores.

